%
% buchcover.tex -- Cover für das Buch Numerik
%
% (c) 2018 Prof Dr Andreas Müller, Hochschule Rapperswil
%
\documentclass[11pt]{standalone}
\usepackage{tikz}
\usepackage{times}
\usepackage{geometry}
\usepackage{german}
\usepackage[utf8]{inputenc}
\usepackage[T1]{fontenc}
\usepackage{times}
\usepackage{amsmath,amscd}
\usepackage{amssymb}
\usepackage{amsfonts}
\usepackage{txfonts}
\usepackage{ifthen}
\usetikzlibrary{math}
\geometry{papersize={402mm,278mm},total={405mm,278mm},top=72.27pt, bottom=0pt, left=72.27pt, right=0pt}
\newboolean{guidelines}
\setboolean{guidelines}{false}
\setboolean{guidelines}{true}

\begin{document}
\begin{tikzpicture}[>=latex, scale=1]
\tikzmath{
	real \ruecken, \einschlag, \gelenk, \breite, \hoehe;
	\ruecken = 3.5;
	\einschlag = 1.6;
	\gelenk = 0.8;
	\breite = 16.7;
	\hoehe = 24.6;
	real \bogengreite, \bogenhoehe;
	\bogenbreite = 2 * (\breite + \einschlag + \gelenk) + \ruecken;
	\bogenhoehe = 2 * \einschlag + \hoehe;
}

%\clip (0,0) circle (6);

\draw[fill=blue](0,0) rectangle({\bogenbreite},{\bogenhoehe});
\hsize=13.6cm

\begin{scope}
\clip (0,0) rectangle({\bogenbreite},{\bogenhoehe});
%\clip (0,0) rectangle ({\bogenbreite},20.2);
%\node at (18.7,8.9) [scale=5.0]{\includegraphics{nozzle-hell.jpg}};
%\node at (18.7,8.9) [scale=5.0]{\includegraphics{nozzle-hell3.jpg}};
%\node at (18.7,8.9) [scale=5.0]{\includegraphics{nozzle.jpg}};
%\node at ({\bogenbreite/2},13.0) {\includegraphics[width=42cm]{matrix.pdf}};
\node at ({\bogenbreite/2},13.0) {\includegraphics[width=42cm]{matrix-blue.png}};
\end{scope}

\node at ({\einschlag+2*\gelenk+\ruecken+1.5*\breite},24.3)
	[color=white,scale=1]
	{\hbox to\hsize{\hfill%
	\sf \fontsize{24}{24}\selectfont Mathematisches Seminar}};

\node at ({\einschlag+2*\gelenk+\ruecken+1.5*\breite},21.9)
	[color=white,scale=1]
	{\hbox to\hsize{\hfill%
	\sf \fontsize{50}{50}\selectfont Matrizen}};

\node at ({\einschlag+2*\gelenk+\ruecken+1.5*\breite},19.7)
	[color=white,scale=1]
	{\hbox to\hsize{\hfill%
	\sf \fontsize{13}{5}\selectfont Andreas Müller und Roy Seitz}};

\node at ({\einschlag+2*\gelenk+\ruecken+1.5*\breite},18.4)
	[color=white,scale=1]
	{\hbox to\hsize{\hfill%
	\sf \fontsize{13}{5}\selectfont
	Joshua Bär,			% E
	Marius Baumann,			% E
	Reto Fritsche,			% E (2)
	Alain Keller%,			% E
%	Ahmet Güzel%,			% E
	}};

\node at ({\einschlag+2*\gelenk+\ruecken+1.5*\breite},17.75)
	[color=white,scale=1]
	{\hbox to\hsize{\hfill%
	\sf \fontsize{13}{5}\selectfont
	Marc Kühne,
	Robine Luchsinger,		% B
	Naoki Pross,			% E
	Thomas Reichlin%,                % B
	}};

\node at ({\einschlag+2*\gelenk+\ruecken+1.5*\breite},17.1)
	[color=white,scale=1]
	{\hbox to\hsize{\hfill%
	\sf \fontsize{13}{5}\selectfont
	Michael Schmid,		% MSE
	Pascal Andreas Schmid,		% B
	Adrian Schuler%,
	}};
 
\node at ({\einschlag+2*\gelenk+\ruecken+1.5*\breite},16.45)
	[color=white,scale=1]
	{\hbox to\hsize{\hfill%
	\sf \fontsize{13}{5}\selectfont
	Thierry Schwaller,		% E
	Michael Steiner,		% E
	Tim Tönz,			% E
	Fabio Viecelli%,		% B
	}};

\node at ({\einschlag+2*\gelenk+\ruecken+1.5*\breite},15.8)
	[color=white,scale=1]
	{\hbox to\hsize{\hfill%
	\sf \fontsize{13}{5}\selectfont
	Lukas Zogg%,			% B
	%
	}};

\node at ({\einschlag+2*\gelenk+\ruecken+1.5*\breite},15.15)
	[color=white,scale=1]
	{\hbox to\hsize{\hfill%
	\sf \fontsize{13}{5}\selectfont
	%Reto Wildhaber%			% B
	}};
 
%\node at (0,3) [color=white] {\sf \LARGE Mathematisches Seminar 2017};

% Rücken
\node at ({\bogenbreite/2 + 0.05},20.5) [color=white,rotate=-90]
	{\sf\fontsize{35}{0}\selectfont Matrizen};

% Buchrückseite
\node at ({\einschlag+0.5*\breite},18.6) [color=white] {\sf
\fontsize{13}{16}\selectfont
\vbox{%
\parindent=0pt
%\raggedright
Das Mathematische Seminar der Ostschweizer Fachhochschule
in Rapperswil hat sich im Frühjahrssemester 2021 dem Thema
Matrizen zugewandt.
Ziel war, die vielfältigen Anwendungsmöglichkeiten dieser grundlegenden
Datenstruktur zu ergründen.
Dieses Buch bringt das Skript des Vorlesungsteils mit den von den
Seminarteilnehmern beigetragenen Seminararbeiten zusammen.
}};


\ifthenelse{\boolean{guidelines}}{
\draw[white] (0,{\einschlag})--({\bogenbreite},{\einschlag});
\draw[white] (0,{\bogenhoehe-\einschlag})--({\bogenbreite},{\bogenhoehe-\einschlag});

\draw[white] ({\einschlag},0)--({\einschlag},{\bogenhoehe});
\draw[white] ({\einschlag+\breite},0)--({\einschlag+\breite},{\bogenhoehe});
\draw[white] ({\einschlag+\breite+\gelenk},0)--({\einschlag+\breite+\gelenk},{\bogenhoehe});
\draw[white] ({\bogenbreite-\einschlag-\breite-\gelenk},0)--({\bogenbreite-\einschlag-\breite-\gelenk},{\bogenhoehe});
\draw[white] ({\bogenbreite-\einschlag-\breite},0)--({\bogenbreite-\einschlag-\breite},{\bogenhoehe});
\draw[white] ({\bogenbreite-\einschlag},0)--({\bogenbreite-\einschlag},{\bogenhoehe});
}{}

\end{tikzpicture}
\end{document}
